(1) 自旋简并度2使动量球内态数 \(N=2V(4\pi k_F^3/3)/(2\pi)^3=Vk_F^3/(3\pi^2)\)，故 \(k_F=(3\pi^2n)^{1/3}\)、\(\varepsilon_F=\hbar^2k_F^2/(2m)\)。三维态密度 \(D(\varepsilon)=C\sqrt\varepsilon\)，\(N=\int_0^{\varepsilon_F}D(\varepsilon)d\varepsilon\)，能量 \(U_0=\int_0^{\varepsilon_F}\varepsilon D(\varepsilon)d\varepsilon=(3/5)N\varepsilon_F\)。由尺度齐次性 \(PV=2U_0/3\)，所以 \(\boxed{P_0=(2/5)n\varepsilon_F}\)。该压力在 \(T=0\) 仍非零，源于Pauli填充而非热运动。

(2) Sommerfeld公式为 \(\int_0^\infty f(\varepsilon)n_F(\varepsilon)d\varepsilon=\int_0^\mu f(\varepsilon)d\varepsilon+(\pi^2/6)(k_BT)^2f'(\mu)+O(T^4)\)。取 \(f=C\sqrt\varepsilon\)，固定 \(N\) 得 \((2C/3)\mu^{3/2}+(\pi^2C/12)(k_BT)^2\mu^{-1/2}=(2C/3)\varepsilon_F^{3/2}\)。设 \(\mu=\varepsilon_F+\delta\mu\) 且 \(\delta\mu=O(T^2)\)，线性展开得 \(C\varepsilon_F^{1/2}\delta\mu+(\pi^2C/12)(k_BT)^2\varepsilon_F^{-1/2}=0\)，故 \(\boxed{\mu(T)=\varepsilon_F[1-(\pi^2/12)(T/T_F)^2+O(T^4/T_F^4)]}\)。

(3) 取能量积分权重 \(f(\varepsilon)=C\varepsilon^{3/2}\)，Sommerfeld展开得 \(U=(2C/5)\mu^{5/2}+(\pi^2C/4)(k_BT)^2\mu^{1/2}+O(T^4)\)。代入上一问 \(\delta\mu\)，整理为 \(\boxed{U/N=(3/5)\varepsilon_F[1+(5\pi^2/12)(T/T_F)^2+O((T/T_F)^4)]}\)。对固定 \(N,V\) 求导，\(\boxed{C_V/Nk_B=(\pi^2/2)(T/T_F)+O((T/T_F)^3)}\)。热容随 \(T\to0\) 线性趋零，符合第三定律；仅有Fermi面附近厚度 \(k_BT\) 的粒子参与低温热激发。

(4) 零温 \(P_0\propto n^{5/3}\)，故 \(\partial P/\partial n=(2/3)\varepsilon_F\)。压缩率 \(\boxed{\kappa_T=[n(\partial P/\partial n)]^{-1}=3/(2n\varepsilon_F)}\)。零温熵为零，等温与绝热微小压缩在此极限一致，声速 \(c_s^2=(1/m)\partial P/\partial n=2\varepsilon_F/(3m)=v_F^2/3\)，即 \(\boxed{c_s=v_F/\sqrt3}\)。若 \(v_F\) 接近光速、相互作用强或系统维度改变，非相对论自由气体的压缩率和声速必须修正。
